% Generated draft: requires mathematical and source review.
\documentclass[11pt]{article}
\usepackage[margin=25mm]{geometry}
\usepackage{fontspec}
\setmainfont{Latin Modern Roman}
\usepackage{amsmath,amssymb,mathtools,unicode-math}
\setmathfont{Latin Modern Math}
\usepackage{enumitem,needspace,xurl,hyperref}
\setlength{\parindent}{0pt}
\setlength{\parskip}{5pt}
\setlength{\emergencystretch}{3em}
\newcommand{\N}{\mathbb{N}}
\newcommand{\Z}{\mathbb{Z}}
\newcommand{\Q}{\mathbb{Q}}
\newcommand{\R}{\mathbb{R}}
\newcommand{\C}{\mathbb{C}}
\newcommand{\F}{\mathbb{F}}
\newcommand{\A}{\mathbb{A}}
\newcommand{\PP}{\mathbb{P}}
\newcommand{\E}{\mathbb{E}}
\newcommand{\eps}{\varepsilon}
\newcommand{\dd}{\,\mathrm{d}}
\begin{document}
\section*{UnsolvedMath ID: 20002787}
\textbf{Source problem code:} catalogue entry
\subsection{Definitions and mathematical statement}
The source asks two related questions: what asymptotic and analytic information about ``basic functions'' is needed to insert them in a trace formula; and how to define Fourier transforms for Vinberg monoids and relate them to Langlands--Shahidi and Rankin--Selberg constructions. No group, monoid, local/global field, representation or trace-formula test space is fixed. The appended ``Endoscopy and Beyond'' is the next heading.

The aid treats only the rank-one toric model $\mathbb G_m\subset\mathbb A^1$. Let $F$ be a nonarchimedean local field, $\mathcal O$ its integers, $\varpi$ a uniformizer, and $q$ its residue cardinality, with $|\varpi|=q^{-1}$. Choose an additive character of conductor $\mathcal O$ and its self-dual additive measure, so $\operatorname{vol}(\mathcal O)=1$. Normalize multiplicative measure by $\operatorname{vol}(\mathcal O^\times)=1$. For $\Phi=\mathbf1_{\mathcal O}$,
\[
 \widehat\Phi=\Phi,\qquad
 \Phi(\varpi^ku)=1\quad(k\ge0,\ u\in\mathcal O^\times).
\]
For an unramified character $\chi:F^\times\to\C^\times$,
\[
 Z(\Phi,\chi,s)=\int_{F^\times}\Phi(x)\chi(x)|x|^s\,\dd^\times x
  =\sum_{k\ge0}(\chi(\varpi)q^{-s})^k
  =L(s,\chi)
\]
initially where $|\chi(\varpi)q^{-s}|<1$. The unramified restriction is essential: with a ramified character this particular unit-invariant test function has zero Mellin integral and is not a universal ramified test vector. The Tate functional-equation convention is
\[
 Z(\widehat\Phi,\chi^{-1},1-s)
       =\gamma(s,\chi,\psi)Z(\Phi,\chi,s).
\]
The aid identifies this abelian factor with the appropriately normalized rank-one Rankin--Selberg/Langlands--Shahidi factors, not with an already constructed transform for every higher monoid.

Its boundary--Mellin dictionary has a precise sequence formulation. For a radial function, or a specified unit-averaged angular component, let $a_k$ be its shell coefficient near zero. Up to a finite Laurent prefix, its unramified Mellin series is $\sum_{k\ge0}a_kz^k$, $z=\chi(\varpi)q^{-s}$. This series is rational exactly when, for all sufficiently large $k$,
\[
 a_k=\sum_\alpha P_\alpha(k)\alpha^k
\]
for finitely many complex $\alpha$ and polynomials $P_\alpha$. A pole at $z=\alpha^{-1}$ of order $m$ corresponds to a nonzero degree-$(m-1)$ polynomial mode, with multiplicities read from the principal parts. Unramified shell averages alone do not determine arbitrary angular behavior. In the basic example, the constant shell mode gives the pole $(1-z)^{-1}$.

For a global field with adeles $\mathbb A$ and a Schwartz--Bruhat function, normalized Poisson summation gives
\[
 \sum_{\alpha\in F}\Phi(\alpha)=\sum_{\alpha\in F}\widehat\Phi(\alpha),
\]
so restriction from the monoid to its nonzero units creates the exact correction
\[
 \sum_{\alpha\in F^\times}\Phi(\alpha)
  -\sum_{\alpha\in F^\times}\widehat\Phi(\alpha)
       =\widehat\Phi(0)-\Phi(0).
\]
Here the global notation is a separate setting from the preceding local field. Mellin pole/boundary-mode control and this missing-orbit correction are independent tests for an intended summation or trace identity. They do not by themselves prove admissibility in a higher-rank trace formula, where the appropriate test space, orbital integrals, convergence or regularization and local normalizations remain unspecified. The source's general Vinberg-monoid questions remain beyond this rank-one audit.
\subsection{Short English statement}
The rank-one toric example links the basic function, Tate Mellin factor and additive Fourier transform. The aid matches Mellin poles to boundary shell modes and shows that deleting the zero orbit adds an explicit Poisson correction; it does not construct the general higher-rank trace-formula theory.
\subsection{Sources}
\begin{itemize}
\item American Institute of Mathematics workshop \emph{Automorphic kernel functions}, first problem group, \url{https://aimath.org/pastworkshops/automorphkernelproblems.pdf}. The requests name Langlands--Shahidi and Rankin--Selberg methods; ``Endoscopy and Beyond'' is an appended next heading.
\item Local UnsolvedMath registry, permanent ID 20002787, source code \texttt{catalogue entry}, canonical location \texttt{aim-representation-theory-notes.json}, \texttt{notes[58]}, record dated 14 August 2026. The original literature field is empty.
\item Supplied research attempt 1: partial result, problem status at run partially solved, rigor labelled proved, novelty confidence low. Its rank-one boundary--Mellin--Poisson audit invokes Tate theory without a full citation and is machine-generated research aid, not peer reviewed; independently verify before citation.
\end{itemize}
\end{document}
